\section{Conclusões}
\label{sec:conclusoes}

Este trabalho estimou como a entrada de data centers e a expansão eólica mudam o balanço horário do subsistema Nordeste em 2030. Os resultados valem para os cenários adotados e respondem à pergunta da Seção~\ref{sec:problema} da seguinte forma.

Com a projeção oficial vigente (\DcRqDois~MWmed) e a expansão eólica base, os data centers aumentam a necessidade de importação em \EfCen~TWh por ano, de \ImpSem{} para \ImpCen~TWh em média, e em \EfHoras{} horas. O aumento equivale a \EfPct\% da energia que consomem. A região continua exportadora, com excedente médio de \ExcCen~TWh, contra \ExcSem~TWh sem data centers. O aumento da importação concentra-se no fim da tarde e à noite.

Compensar a energia anual dos data centers exige \CompEnergia~GW de eólica adicional. Compensar o aumento da importação exige \CompImp~GW, porque a eólica adicional produz pouco nas horas em que falta geração. Um balanço anual equilibrado não garante, portanto, o atendimento horário.

A importação varia muito com o clima do ano: no cenário central, o intervalo de P5 a P95 vai de \ImpCenPCinco{} a \ImpCenPNoventa~TWh. Usar 31 anos de reanálise, em vez dos quatro anos com dados do ONS, foi necessário para representar essa variação.

Do ponto de vista dos dados, três decisões foram determinantes: usar uma série de carga com definição homogênea, modelar o CF eólico potencial no lugar do observado e escolher as variáveis por ablação fora da amostra. As verificações mostraram também o principal limite do método, que é subestimar a importação fora da amostra, sobretudo em 2025 e 2026. O efeito dos data centers varia menos que o nível entre as hipóteses testadas, de \EfSensMin{} a \EfSensMax~TWh, exceto com atraso na sua entrada.

As demais limitações são a ausência da rede elétrica, a hidráulica representada por uma regra, a térmica fixada no histórico, os quatro anos hidrológicos e a carga constante dos data centers. Os próximos passos são representar o redespacho e os limites de intercâmbio, ampliar a série hidrológica e avaliar o armazenamento e a flexibilidade da carga como alternativas à expansão eólica.
